\section{Scope and obstructions}\label{sec:limitations}

The aggregate argument controls the sum of squared copy errors. It does not
assert that every bag copy is within a constant number of core widths of
the optimizer. We first show that such a uniform localization assertion
fails on branching trees, even when all structural and conditioning
constants are fixed. We then show why a geometric feasible repair alone
does not extend the objective-gradient slope rule to coupled constraints.

\subsection{Uniform localization fails on branching trees}
\label{sec:limitations-localization}

A possible alternative to aggregate contraction would bound every copy of
every minimizing configuration by $KW$ whenever all selected boxes have
width at most $W$. The following example rules out that assertion with
$K$ depending only on occurrence, width, branching, and conditioning.
The configuration minimum is the one in \cref{prop:model-dp}, with all
nonempty closed intersections retained.

\begin{theorem}[Failure of uniform copy localization]
\label{thm:limitations-localization}
There is a family satisfying \cref{eq:model-contract,eq:model-growth}
with $p=2$, $q=1$, $k=3$, at most two children per bag, $M=3$, $A=1/2$,
and a fixed $g>0$, whose optimizer is $x^*=0$ and whose exact
optimizer-based separator slopes are zero, with the following property.
For every $K>0$, one member admits finite dyadic bag and separator
partitions of width at most $W$ such that every minimizing configuration
has its root coordinate greater than $KW$. The objective is a strongly
convex quadratic with a spectral condition number bounded independently
of the number of bags.
\end{theorem}

\begin{proof}
Fix $a=7/20$ and let $\mathcal G_d$ be a complete binary tree of depth
$d\ge1$, with $n=2^{d+1}-1$ vertices, root $o$, adjacency matrix
$\mathsf A_d$, and graph-leaf set $\mathcal V_d$. Set
\[
 H=I-a\mathsf A_d,\qquad \delta=1-2\sqrt2a>0,
 \qquad
 F(x,y)=\frac12x^{\mathsf T}Hx+
              \frac12\sum_{v\in\mathcal V_d}y_v^2
 \quad\text{on }[-1,1]^{n+|\mathcal V_d|}.
\]
To bound the adjacency norm, orient each edge from its parent to its
child and apply
$2|uv|\le u^2/\sqrt2+\sqrt2v^2$.
Each vertex is a parent on at most two edges and a child on at most one,
so $|x^{\mathsf T}\mathsf A_dx|\le2\sqrt2\norm{x}_2^2$. Thus
\begin{equation}\label{eq:limitations-spectrum}
 \delta I\preceq H\preceq(1+2\sqrt2a)I.
\end{equation}
The unique optimizer is zero and $g=\delta/2$ is valid in every
dimension. The full objective Hessian has condition number at most
$(1+2\sqrt2a)/\delta$.

Use a root bag $\{x_o\}$. For each nonroot graph vertex $v$, use the
bag $\{x_{\operatorname{par}(v)},x_v\}$ and attach it to the bag of its
graph parent. At every graph leaf $v$, attach a final bag $\{x_v,y_v\}$.
The running-intersection property holds. Internal graph coordinates
occur in their own bag and their two children's bags; graph-leaf
coordinates occur in their own bag and the final bag. Hence $k=3$,
$p=2$, $q=1$, and branching is at most two.

We next specify a factorization and its local models. Put
\[
 \sigma=\sqrt{1-8a^2},\qquad \lambda_+=(1+\sigma)/2,
 \qquad
 P_v=\begin{cases}
  2,&v\in\mathcal V_d,\\
  1-\displaystyle\sum_{u\text{ child of }v}a^2/P_u,&v\notin\mathcal V_d.
 \end{cases}
\]
All pivots lie in $[\lambda_+,2]$: the map
$P\mapsto1-2a^2/P$ preserves that interval and fixes $\lambda_+$.
Expanding squares gives the identity
\begin{equation}\label{eq:limitations-factorization}
 F(x,y)=\frac{P_o}2x_o^2+
 \sum_{v\ne o}\frac{P_v}2
       \left(x_v-\frac a{P_v}x_{\operatorname{par}(v)}\right)^2
 +\sum_{v\in\mathcal V_d}\left(-\frac{x_v^2}2+\frac{y_v^2}2\right).
\end{equation}
Indeed, every off-diagonal coefficient is $-a$, every nonleaf diagonal
coefficient is $(P_v+\sum_ua^2/P_u)/2=1/2$, and every graph-leaf
diagonal coefficient is $(2-1)/2=1/2$.
Assign each displayed summand to its corresponding bag. Keep the root
and edge squares exact, and keep $y_v^2/2$ exact. On an interval $[l,u]$,
replace $-z^2/2$ by its chord
\[
 -\frac{z^2}2-\frac12(z-l)(u-z).
\]
This is an affine valid lower model. Its error is at most
$(u-l)^2/8$, so the bag contract holds with $A=1/2$. The edge-square
Hessian norm is $P_v+a^2/P_v<3$; the root norm is at most $2$ and the
final-bag norm is $1$. Thus $M=3$ is valid. Every assigned bag gradient
vanishes at zero, so all slopes in \cref{eq:model-slope} at zero vanish.

For a dyadic $0<W\le1$, define
\[
 h_W(z)=\begin{cases}z(W-z),&0\le z\le W,\\0,&\text{otherwise},\end{cases}
 \qquad
 G_W(x,y)=F(x,y)-\frac12\sum_{v\in\mathcal V_d}h_W(x_v).
\]
We compute its unique global minimizer before constructing the finite
partitions. Partition the graph coordinates into interior vertices $I$
and leaves, write $z=x_{\mathcal V_d}$, and set
\[
 B=-H_{II}^{-1}H_{I\mathcal V_d},\qquad
 S=H_{\mathcal V_d\mathcal V_d}
       -H_{\mathcal V_d I}H_{II}^{-1}H_{I\mathcal V_d}.
\]
Completing the square yields
\begin{equation}\label{eq:limitations-schur}
 G_W=\frac12(x_I-Bz)^{\mathsf T}H_{II}(x_I-Bz)
       +\frac12z^{\mathsf T}Sz-\frac12\sum_vh_W(z_v)
       +\frac12\norm{y}_2^2.
\end{equation}
The convergent Neumann series for $H_{II}^{-1}$ has nonnegative
entries. Thus $S$ has nonpositive off-diagonal entries, and symmetry
makes its row sums equal, say to $q_d$.

To compute $q_d$, harmonically extend the constant leaf vector $1$.
With
\[
 r_\pm=\frac{1\pm\sigma}{4a},\qquad \gamma=r_-/r_+,
\]
its value at graph level $j$ is
\[
 v_j=\frac{r_+^{j+1}-r_-^{j+1}}
              {r_+^{d+1}-r_-^{d+1}}.
\]
The roots satisfy $2ar^2-r+a=0$, so
$v_0=2av_1$, $v_j=av_{j-1}+2av_{j+1}$ for $0<j<d$, and $v_d=1$.
These are exactly the interior equations, whose solution is unique by
\cref{eq:limitations-spectrum}. The residual in a leaf row is therefore
\[
 q_d=1-av_{d-1}
     =\lambda_+\frac{1-\gamma^{d+2}}{1-\gamma^{d+1}},
 \qquad \lambda_+<q_d\le1.
\]
The last upper bound also follows from $v_{d-1}>0$.

Write $b_{ij}=-S_{ij}\ge0$ for distinct leaves. The reduced part of
\cref{eq:limitations-schur} is
\[
 \sum_i\psi(z_i)+\frac12\sum_{i<j}b_{ij}(z_i-z_j)^2,
 \qquad \psi(z)=\frac{q_d}2z^2-\frac12h_W(z).
\]
Its scalar minimizer is $t_d=W/[2(q_d+1)]$. Inside $[0,W]$,
$\psi(z)-\psi(t_d)=(q_d+1)(z-t_d)^2/2$. Outside that interval,
\[
 \psi(z)-\psi(t_d)-\frac{q_d}4(z-t_d)^2
       =\frac{q_d}4(z+t_d)^2+\frac{t_d^2}2\ge0.
\]
Consequently the unrestricted auxiliary minimizer is unique: every
graph leaf equals $t_d$, $y=0$, and the interior vector is its harmonic
extension. Its coordinate at level $j$ is
\[
 \bar x_j=t_d
       \frac{r_+^{j+1}-r_-^{j+1}}{r_+^{d+1}-r_-^{d+1}}.
\]
Since $a>1/3$, $r_+<1$, and hence
\begin{equation}\label{eq:limitations-amplification}
 \frac{\bar x_0}{W}
   =\frac{(1-\gamma)r_+^{-d}}
           {2(q_d+1)(1-\gamma^{d+1})}\longrightarrow\infty.
\end{equation}
At each depth choose $W$ dyadic and small enough that every $\bar x_j$
lies strictly inside $[-1/2,1/2]$. The ratio in
\cref{eq:limitations-amplification} is unchanged, and this unrestricted
minimizer is also the unique minimizer on the domain box.

We need a growth bound for the auxiliary objective. By
\cref{eq:limitations-spectrum},
$\norm{B}_2\le b_0:=2\sqrt2a/\delta$.
The preceding scalar inequality and \cref{eq:limitations-schur} show
\[
 E:=G_W(x,y)-G_W(\bar x,0)
 \ge\frac\delta2\norm{x_I-Bz}_2^2
       +\frac{q_d}4\norm{z-t_d\one}_2^2+\frac12\norm{y}_2^2.
\]
Using $\bar x_I=B(t_d\one)$ and
\[
 \norm{x_I-\bar x_I}_2^2
 \le2\norm{x_I-Bz}_2^2+2b_0^2\norm{z-t_d\one}_2^2,
\]
we obtain, for the depth-independent constant
\[
 C_*:=\frac4\delta+\frac{4(2b_0^2+1)}{\lambda_+},
 \qquad
 \norm{(x,y)-(\bar x,0)}_2^2\le C_*E.
\]
Here $C_*>2$, so it also covers the $y$ coordinates.

Choose a dyadic $0<h\le W$ dividing $W$. Partition all original root
and edge bags into uniform cubes of width $h$, and every separator into
intervals of width $h$. In each final bag $\{x_v,y_v\}$, use squares of
width $W$ on the slab $[0,W]\times[-1,1]$ and squares of width $h$ on
its complement. These are finite dyadic partitions: start with the
width-$W$ grid and refine the squares outside the slab. Every leaf and
cell has width at most $W$. Use zero slopes and the maximal finite
intercepts of \cref{prop:model-dp}.

For any configuration $c$, form its top-copy point $(x(c),y(c))$.
The parent coordinate in an original edge bag and the graph-leaf
coordinate in a final bag differ from their top copies by at most $2h$:
the relevant separator interval has width $h$ and meets a parent
interval of width $h$. All other coordinates of those bags already
are top copies. On $[-1,1]^2$, the gradient norm of an edge square is
at most $3\sqrt2$, so changing its parent copy alters its value by
at most $12h$.

Put $\varphi_W(z)=-z^2/2-h_W(z)/2$. It is $1$-Lipschitz on
$[-1,1]$: its derivative is $-z$ outside $[0,W]$ and $-W/2$ inside.
On a coarse final square, the chord model is exactly $\varphi_W$.
On a fine final square it is $\varphi_W-e$, with $0\le e\le h^2/8$.
Its discrepancy from evaluation at the top copy is thus at most
$2h+h^2/8$. The root and dummy-coordinate factors require no change.
Summing, since $h\le1$ and $|\mathcal V_d|\le n$, gives the uniform
configuration estimate
\begin{equation}\label{eq:limitations-finite-realization}
 \bigl|\Phi(c)-G_W(x(c),y(c))\bigr|\le15nh.
\end{equation}
This estimate includes every configuration using boundary-only
intersections.

There is also a consistent configuration at $(\bar x,0)$. Its graph
leaves lie strictly in $[0,W]$, so choose their coarse final squares.
All convex factors are exact and the concave chords supply precisely
the auxiliary wells. Hence $\min_c\Phi(c)\le G_W(\bar x,0)$.
For every minimizing configuration, \cref{eq:limitations-finite-realization}
and auxiliary growth imply
\[
 \norm{(x(c),y(c))-(\bar x,0)}_2^2\le15C_*nh.
\]
Given $K$, first choose $d$ with $\bar x_0/W>2K$, then choose the
domain-compatible dyadic $W$, and finally choose $h$ dividing $W$ so
small that $15C_*nh<\bar x_0^2/4$. Every minimizing configuration then
has root coordinate greater than $\bar x_0/2>KW$, as required.
\end{proof}

The example satisfies the stronger grading condition
$\width(Y)\le W+\theta\dist_\infty(Y,x^*)$ with $\theta=0$, for
every leaf and cell. Exact slopes therefore do not repair unrestricted
tree localization. The true objective is strongly convex, but its supplied
factorization contains concave leaf factors. Keeping another convex
factorization exact would change the local relaxation.

This theorem concerns arbitrary allowed partitions. It does not show that
a specified nested-refinement algorithm reaches these partitions, give a
lower bound on the smallest certificate, or preclude a localization theorem
under an additional influence-contraction assumption. It also does not
contradict aggregate contraction: a root error much larger than $W$ can
still be compatible with a bound on total squared error that depends on the
number of bags.

\subsection{Feasible repair does not cancel first-order error}
\label{sec:limitations-constraint}

\begin{proposition}[An affine-constraint obstruction]
\label{prop:limitations-constraint}
Let $L>0$ and minimize
\[
 F(x,s,y)=\frac{x^2+s^2+y^2}{2}+Lx-Ly
 \quad\text{on }[-1,1]^3,\qquad x=s,\quad y=s.
\]
Use root bag $\{x,s\}$ with objective
$a_r=(x^2+s^2)/2+Lx$ and child bag $\{s,y\}$ with objective
$a_t=y^2/2-Ly$, assigning each equality to its corresponding bag.
Take exact local objective models and enforce the local equalities in
the convex subproblems. The feasible problem has point growth $g=1/2$,
bag curvature $M=1$, zero objective-model error, and a box-preserving
feasible repair with Euclidean Lipschitz constant one. Nevertheless,
any valid certificate with constant child cell minorants proving
tolerance $0<\eps\le3/2$ requires
\begin{equation}\label{eq:limitations-constraint-count}
 |\mathcal P_t|\ge\sqrt{\frac23}\frac L{\sqrt\eps}.
\end{equation}
The unmodified objective-gradient slope is zero at every center. A
multiplier-adjusted slope gives an exact certificate with one leaf per
bag and one separator cell.
\end{proposition}

\begin{proof}
The feasible set is $\{(v,v,v):-1\le v\le1\}$, on which
$F=3v^2/2=\norm{(v,v,v)}_2^2/2$. Thus zero is the unique optimizer
and $g=1/2$. Both bag Hessian norms are at most one. The orthogonal
projection
\[
 R(z)=\frac{z_x+z_s+z_y}{3}(1,1,1)
\]
maps the box into the feasible set and is $1$-Lipschitz. It also obeys
\[
 \norm{z-R(z)}_2^2
 =\frac{(z_x-z_s)^2+(z_s-z_y)^2+(z_x-z_y)^2}{3}
 \le(z_x-z_s)^2+(z_y-z_s)^2.
\]
In particular even residual-to-repair distance has a fixed constant.
The prescribed child slope is $\partial_s a_t=0$, independently of
the center.

Let $D$ be any separator interval and let $u\le v$ belong to $D$.
The child certificate inequality at its locally feasible point $(v,v)$
forces its constant intercept to satisfy
$\beta_D\le v^2/2-Lv$. At the root point $(u,u)$, choose any leaf
containing that point. Its convex child bound is at most $\beta_D$
at $u$ by the child-minorant test. Consequently
\begin{equation}\label{eq:limitations-constraint-pair}
 \LB\le u^2+Lu+\beta_D
       \le u^2+v^2/2-L(v-u).
\end{equation}
This argument permits arbitrary finite bag partitions and arbitrary
valid convex child bounds. Weaker objective models can only lower the
right-hand comparison.

If the certificate proves tolerance $\eps$, then $\LB\ge-\eps$.
Fix $0<r\le1$, intersect each cell with $[-r,r]$, and use the
endpoints of that intersection for $u,v$ in
\cref{eq:limitations-constraint-pair}. It follows that
\[
 L\,\operatorname{length}(D\cap[-r,r])\le\eps+3r^2/2.
\]
The cell interiors are disjoint and cover the separator domain, so
their intersection lengths sum to $2r$. Thus
$|\mathcal P_t|\ge2rL/(\eps+3r^2/2)$. Choosing
$r=\sqrt{2\eps/3}\le1$ proves
\cref{eq:limitations-constraint-count}.

For comparison, take the single cell $[-1,1]$ with minorant
$\ell_t(s)=-Ls$. It is below the exact child objective on $y=s$,
because their difference is $s^2/2$. Adding it at the root on $x=s$
gives $s^2\ge0$. Thus the resulting root bound is exactly zero.
The equality multipliers for the convention
$F+\mu_1(x-s)+\mu_2(y-s)$ at zero are $\mu_1=-L$ and
$\mu_2=L$; the adjusted child separator derivative is $-L$.
\end{proof}

For example, a cell $[0,h]$ admits copies $(0,0)$ and $(h,h)$ whose
configuration value is $h^2/2-Lh$. Repairing the top-copy point
$(0,0,h)$ produces objective $h^2/6$, so its difference from that
configuration value is $Lh-h^2/3$. The first-order coefficient is not
controlled by objective curvature or a geometric repair constant.
Constraint-aware slopes must supply the cancellation used by the dynamics
construction. In nonlinear dynamics, bounding the resulting multipliers
also requires control of objective first derivatives.

\subsection{Parameters, output, and other optimal sets}
\label{sec:limitations-parameters}

The point-growth assumption \cref{eq:model-growth} forces a unique
minimizer: applying it at any other minimizer gives zero distance from
$x^*$. Growth toward an entire optimal set is a different hypothesis.
The contraction proof here does not discover or retain the multiple
centers needed to handle separated optimal components. An extension could
retain several centers or exploit a description of optimal coordinate
projections, but it would require a new progress argument and a bound on
the additional representation.
The lower allocation result in \cref{sec:allocation} requires an additional
lower bound on local model error; it does not supply a matching upper
construction for general degenerate optimal sets.

Occurrence and conditioning enter the construction through the supplied
bag functions. Neither treewidth nor a bound on the full objective Hessian
controls the number of occurrences or the bag-gradient constants. Changing
factor assignments, adding cancelling bag terms, or retaining useful convex
terms exactly can change the local models and their errors. Thus the
parameters describe the supplied representation. They are not invariant
properties of the objective alone. A small growth constant caused by a
distant nearly optimal point also remains a substantive limitation; the
algorithm does not verify the global growth promise.

Fresh regridding produces a new certificate at each stage. Its final
member count, total number of created boxes, and local convex-oracle count
are different quantities. Internal oracle work, factor evaluation costs,
and the length of serialized local proofs require their own representation
and precision bounds. The arithmetic specializations establish such bounds
only for their stated data and output contracts. In particular the nonlinear
output is an exact recurrence from rational initial state and controls,
together with certified objective bounds; it need not list expanded rational
state coordinates.

Unknown-conditioning searches preserve validity by stopping only at a
computed certificate gap and enforcing budgets before work is performed.
Their work guarantee compares total search with a sufficient trial. The
first successful trial may use a larger grading ratio and reach a later
stage than that sufficient trial. The sharp final-size formula for a
sufficient fixed-ratio run therefore must be distinguished from the total
work and output-budget guarantee of the search. A box-only budget does not
by itself bound incidence processing, oracle-internal work, or bit
operations.

The general result is for product boxes. The dynamics extension uses a
specific box-invariant stable simulation and valid local affine outer
strips, with no extra terminal or state constraints. Arbitrary coupled
constraints, merely approximate local feasibility, or a repair controlled
only in distance are not covered. A curvature bound used to build a graph
enclosure is part of certificate soundness; decreasing the grading ratio
cannot make an underestimated enclosure valid. These restrictions separate
the validity of an emitted certificate from the assumptions that give the
stated rate and computational bound.
